\section{System Model and Governing Principles}
\label{sec:circuit}

\subsection{Circuit Overview}

We consider a series RLC circuit driven by a step voltage source $V(t)$.
At $t = 0$, a DC source of amplitude $\Vs$ is connected to an initially
uncharged capacitor, so $V(t) = \Vs$ for the post-switching interval
$t \geq 0$.
The quantity of primary interest is the capacitor voltage $\vC(t)$ for
$t \geq 0$.
The component values used for the main RLC case study are summarized in
Table~\ref{tab:nominal-values}.

\begin{table}[!t]
  \caption{Nominal Component Values}
  \label{tab:nominal-values}
  \centering
  \begin{tabular}{@{}lll@{}}
    \toprule
    Quantity & Symbol & Nominal value \\
    \midrule
    Source amplitude & $\Vs$ & $10\,\mathrm{V}$ \\
    Resistance cases & $R$ & $50\,\Omega$, $200\,\Omega$, $500\,\Omega$ \\
    Inductance & $L$ & $100\,\mathrm{mH} = 0.1\,\mathrm{H}$ \\
    Capacitance & $C$ & $10\,\mu\mathrm{F} = 10\times 10^{-6}\,\mathrm{F}$ \\
    \bottomrule
  \end{tabular}
\end{table}

These values were chosen because they produce simple, transparent damping
parameters while remaining realistic for a bench-top laboratory circuit:
$LC = 1\times 10^{-6}\,\mathrm{s}^2$, so
$\omega_0 = 1/\sqrt{LC} = 1000\,\mathrm{rad/s}$, and
$\Rcrit = 2\sqrt{L/C} = 200\,\Omega$.
Thus $R = 50\,\Omega$, $R = 200\,\Omega$, and
$R = 500\,\Omega$ provide clear underdamped, critically damped, and
overdamped examples, respectively.

\subsection{State Variables}

The two state variables are the capacitor voltage $\vC(t)$ and the inductor
current $i(t)$, which also equals the series current through every element.
These quantities are chosen because neither can change instantaneously:
the capacitor voltage is proportional to stored charge ($\vC = q/C$), and the
inductor current is proportional to stored magnetic flux ($\Psi = Li$).
Consequently, they fully characterize the energy state of the circuit at any
instant.

\subsection{Initial Conditions}

Before the switch is closed at $t = 0$:
\begin{itemize}
  \item The capacitor is initially uncharged: $\vC(0) = 0\,\mathrm{V}$.
  \item No current flows: $i(0) = 0\,\mathrm{A}$.
\end{itemize}
Prior to switching, the capacitor behaves as an open circuit (blocking DC
current) and the ideal inductor behaves as a short circuit (zero DC voltage
drop).

\subsection{Governing Equations}

\subsubsection{Kirchhoff's Voltage Law}

Applying Kirchhoff's Voltage Law (KVL) around the series loop~\cite{nilsson2019}:
\begin{equation}
  Ri + L\frac{di}{dt} + \vC = V(t).
  \label{eq:kvl}
\end{equation}

\subsubsection{Second-Order ODE for the RLC Circuit}

Using the capacitor voltage relationship $i(t) = C\,d\vC/dt$ and
substituting into \eqref{eq:kvl} yields the governing second-order ODE:
\begin{equation}
  LC\frac{d^2\vC}{dt^2} + RC\frac{d\vC}{dt} + \vC = V(t).
  \label{eq:rlc2nd}
\end{equation}

\subsubsection{First-Order ODEs for Reduced Circuits}

For an RC circuit (inductor replaced by an ideal wire, $L \to 0$):
\begin{equation}
  \frac{d\vC}{dt} + \frac{1}{RC}\,\vC = \frac{\Vs}{RC}.
  \label{eq:rc1st}
\end{equation}

For an RL circuit (capacitor replaced by an ideal wire):
\begin{equation}
  \frac{di}{dt} + \frac{R}{L}\,i = \frac{\Vs}{L}.
  \label{eq:rl1st}
\end{equation}
